%=====================================================================
% LEITURA E DIAGNOSTICO POR ESQUEMAS
%=====================================================================
\section{Leitura, identificação e diagnóstico}

\begin{frame}{Da prancha à instalação — sequência única de leitura}
\centering
\begin{tikzpicture}[
  etapa/.style={draw=corPrim,fill=corDef!5,rounded corners,minimum width=2.05cm,minimum height=.82cm,align=center,font=\scriptsize},
  fluxo/.style={-{Stealth},thick,corPrim}]
\node[etapa](p)at(-5.2,0){planta\\onde está};
\node[etapa](q)at(-2.6,0){quadro\\de onde vem};
\node[etapa](d)at(0,0){diagrama\\como funciona};
\node[etapa](b)at(2.6,0){bornes\\como conecta};
\node[etapa](e)at(5.2,0){ensaio\\como verifica};
\draw[fluxo](p)--(q)--(d)--(b)--(e);
\end{tikzpicture}
\vspace{4mm}
\begin{caixaProjeto}[title=Critério de interpretação]
O aluno deve conseguir apontar a mesma função em todas as representações: ponto na planta, circuito no quadro, símbolo no diagrama, borne no componente e item do ensaio.
\end{caixaProjeto}
\end{frame}

\begin{frame}{Numeração de fios e bornes — transformar desenho em conexão rastreável}
\small
\begin{columns}[T,onlytextwidth]
\column{.56\textwidth}
\centering
\begin{tikzpicture}[scale=.82,every node/.style={font=\scriptsize}]
\node[draw=corPrim,minimum width=1.8cm](x1)at(-3.1,1.5){X1:01};
\node[draw=corNorma,minimum width=1.8cm](s1)at(0,1.5){S1:L};
\node[draw=corNorma,minimum width=1.8cm](s2)at(0,-.2){S1:1};
\node[draw=corPrim,minimum width=1.8cm](x2)at(3.1,-.2){X2:07};
\draw[thick,corPrim,-{Stealth}](x1)--(s1)node[midway,above]{101};
\draw[thick,corAccent,-{Stealth}](s2)--(x2)node[midway,above]{102};
\draw[dashed,corPrim!60](s1)--(s2);
\node[font=\tiny]at(.55,.65){contato};
\end{tikzpicture}
\column{.42\textwidth}
\begin{itemize}
\item X1/X2: régua ou caixa;
\item 01/07: posição física;
\item S1:L e S1:1: bornes do dispositivo;
\item 101/102: identificação do condutor;
\item a mesma referência deve aparecer na documentação de campo.
\end{itemize}
\end{columns}
\end{frame}

\begin{frame}{Como interpretar um esquema desconhecido}
\small
\begin{enumerate}
\item localizar título, revisão e tensão indicada;
\item marcar origem, proteção e referências de potencial;
\item separar potência, comando, sinal e proteção;
\item seguir um caminho por vez, sem saltar entre ramos;
\item identificar dispositivos pelos seus bornes e contatos;
\item registrar origem e destino de cada fio/cabo;
\item comparar a representação com a marcação do componente;
\item conferir se planta, quadro, diagrama e lista de cabos contam a mesma história.
\end{enumerate}
\begin{caixaAt}
Se um condutor não tiver função e destino claros no desenho, o esquema ainda não está suficientemente compreendido para servir de guia de montagem.
\end{caixaAt}
\end{frame}

\begin{frame}{Diagnóstico por nós — evitar tentativa e erro}
\centering
\begin{tikzpicture}[
  n/.style={draw=corPrim,fill=corDef!5,rounded corners,minimum width=2.15cm,minimum height=.78cm,align=center,font=\scriptsize},
  a/.style={-{Stealth},thick,corPrim}]
\node[n](doc)at(-5.0,1.1){1. conferir\\documentação};
\node[n](vis)at(-2.5,1.1){2. inspeção\\visual};
\node[n](orig)at(0,1.1){3. verificar\\origem};
\node[n](traj)at(2.5,1.1){4. rastrear\\trajeto};
\node[n](carga)at(5.0,1.1){5. verificar\\destino};
\draw[a](doc)--(vis)--(orig)--(traj)--(carga);
\node[n,draw=corNorma,fill=corNorma!5,minimum width=3.3cm](reg)at(0,-1.15){registrar causa, correção\\e nova verificação};
\draw[a](carga.south)--(5.0,-1.15)--(reg.east);
\end{tikzpicture}
\vspace{2mm}
\begin{caixaObs}
O diagrama deve permitir localizar a falha por etapas. Desenhos genéricos demais obrigam o aluno a adivinhar onde começa e termina cada função.
\end{caixaObs}
\end{frame}
